\documentclass[tikz,border=8pt]{standalone}
\usepackage{ctex}
\usepackage{amsmath}
\usepackage{amssymb}
\usepackage{pgfplots}
\pgfplotsset{compat=1.18}
\usetikzlibrary{calc,angles,quotes,arrows.meta,positioning,decorations.markings}
\begin{document}
\begin{tikzpicture}[x=1.0cm, y=1cm]
% ---- 上：原始线性尺度的困境（示意）----
\node[anchor=west, font=\small, black!70] at (-0.3,3.1)
  {线性尺度上 $1$ 与 $10^{14}$ 相差 14 个数量级，无法同框显示};

% ---- 主：对数标度数轴 ----
% 把 10^0 .. 10^14 等距放到 0..14 (每格 1cm)
\draw[-{Latex[length=3mm]}, very thick] (-0.5,0) -- (15.2,0);
\node[anchor=west, font=\small] at (15.0,-0.55) {$\lg x$ 轴（等距）};

\foreach \k in {0,1,2,3,4,5,6,7,8,9,10,11,12,13,14}{
  \draw[thick] (\k,0.18) -- (\k,-0.18);
}
% 主刻度标注 10^k
\foreach \k in {0,2,4,6,8,10,12,14}{
  \node[anchor=north, font=\footnotesize] at (\k,-0.22) {$10^{\k}$};
}

% 顶部说明框（半透明，置于上方空白处，避开刻度）
\node[fill=yellow!16, draw=black!35, rounded corners, align=center,
      fill opacity=0.8, text opacity=1, font=\small]
      at (7,2.25)
      {对数标度：把跨 14 个数量级的量 $1,10,10^2,\dots,10^{14}$\\
       压缩到\textbf{等距}的线性轴上（相邻刻度均为 10 倍）};

% ---- 现实实例标注（下方，彩色指引线，错落避免重叠）----
% pH：氢离子浓度 10^0 (酸) 到约 10^{-14}? 这里以"浓度比"正向放置示意：差1个pH=10倍
\draw[blue!70!black, thick] (1,0) -- (1,-1.0);
\node[blue!70!black, anchor=north, font=\footnotesize, align=center] at (1,-1.05)
  {pH\\差 $1$ 单位\\$=$ 浓度 $10$ 倍};

\draw[red!75!black, thick] (5,0) -- (5,-1.6);
\node[red!75!black, anchor=north, font=\footnotesize, align=center] at (5,-1.65)
  {分贝 dB\\每 $+10$ dB\\$=$ 声强 $\times 10$};

\draw[green!45!black, thick] (9,0) -- (9,-1.0);
\node[green!45!black, anchor=north, font=\footnotesize, align=center] at (9,-1.05)
  {里氏震级\\震级 $+1$\\$=$ 振幅 $\times 10$};

\draw[purple!75!black, thick] (13,0) -- (13,-1.6);
\node[purple!75!black, anchor=north, font=\footnotesize, align=center] at (13,-1.65)
  {天文亮度\\跨度极大\\需压缩尺度};

\end{tikzpicture}
\end{document}
